\section{相关工作}

\subsection{遥感图文表征与开放词汇分割}
CLIP 通过图文对比学习建立视觉与语言的对应关系，为文本指定类别提供表征基础~\cite{radfordLearningTransferableVisual2021}。RemoteCLIP 将遥感检测框和分割掩码转换为图文样本，扩展遥感领域的图文预训练资源~\cite{liuRemoteCLIPVisionLanguage2024}。其领域表征能够提供语义先验，而将这类先验转化为像素级预测仍需要匹配、空间建模与解码设计。

GSNet 通过类别查询引导通用 CLIP 与遥感专用视觉特征融合，并以残差信息保留解码器恢复细节；其原始分割适配使用 LandDiscover50K~\cite{yeOpenVocabularyRemoteSensing2025}。RSKT-Seg 聚合多方向图文匹配，并通过 RemoteCLIP 等预训练编码器的中间特征进行领域知识迁移~\cite{liExploringEfficientOpenVocabulary2026}。Pi-Seg 在文本原型和稠密视觉特征中引入可学习扰动，以分割监督约束图文对齐~\cite{liPerturbationRegularizedOpenVocabularyRemote2026}。SegEarth-OV 则利用无标注遥感图像训练特征上采样器，并在推理时削弱局部特征中的全局偏置~\cite{liSegEarthOVTrainingFreeOpenVocabulary2025}。这些方法分别从表征、对齐或空间恢复入手。本文将遥感相似度编码为类别相关指导特征，使其参与 SAM 3 条件特征的注意力更新。

\subsection{SAM 3 与遥感提示适配}
SAM 3 支持以短名词短语或视觉示例指定待分割概念，并提供实例掩码及概念相关的语义响应~\cite{carionSAM3Segment2026}。SegEarth-OV3 逐类别查询 SAM 3，融合实例与语义响应，并使用概念存在分数抑制缺席类别~\cite{liSegEarthOV3ExploringSAM2026}。ProC-SAM3 构建离线类别提示池、缓存文本嵌入，并通过存在性和置信度校准融合响应~\cite{songPromptCalibratedSAM32026}。ROSE-SAM3 利用无标签图像上的冻结响应离线选择提示，并在预测时进行语义与实例协作~\cite{songROSESAM3ResponseOrientedPrompt2026}。RSGPNet 将初始预测掩码转换为几何框提示，再通过一致性验证细化结果~\cite{wangRSGPNetGeometricPrompting2026}。这些方法在保持模型参数冻结的同时调整提示或推理流程。ActiveSAM 则根据图像中的类别存在性裁剪活动词汇以减少解码开销~\cite{tienActiveSAMFastAccurate2026}。

本文使用 DLRSD 分割监督学习类别条件特征聚合，并在固定权重下进行目标域推理。

上述路径主要调整提示或预测输出。本文在解码前引入可学习的跨提示特征交互，采用源域监督训练聚合和融合模块，并部分适配 RemoteCLIP。SAM 3 参数保持冻结，原始语义路径仅在训练时提供软掩码约束；推理使用学生语义路径。参考掩码为特征适配提供原始空间响应的一致性约束。

\subsection{空间、类别与结构引导聚合}
CAT-Seg 将图文余弦相似度组织为代价体，交替进行空间和类别聚合，并通过上采样解码生成预测~\cite{choCATSegCostAggregation2024}。PCA-Seg 将两类聚合并行处理后融合~\cite{yinPCASegRevisitingCost2026}。RSKT-Seg 进一步设计空间及类别维度的高效代价融合~\cite{liExploringEfficientOpenVocabulary2026}。GeoSeg-OV 利用冻结视觉基础模型的结构特征，引导代价传播和多尺度解码，体现了语义匹配与空间结构协同建模的另一种方式~\cite{liuGeoSegOV2026}。

这些工作已经研究了空间与类别上下文，本文承接这一思想并改变交互对象与解码接口：模板级相似度经卷积和最终稠密 RemoteCLIP 特征融合后，形成每个提示的语义引导；它与当前 SAM 3 类别条件特征共同构造注意力的查询和键，而值始终来自条件特征。每个聚合块依次进行类间交互与共享权重的多尺度窗口聚合，仅残差更新条件特征。语义引导在各块间保持不变，并通过查询和键参与训练。最终送入冻结分割解码器的是聚合后的 SAM 3 特征经 SDF 融合得到的表征，相似度由此引导条件分割特征的适配。
